\section{Einführung zur KI}
\subsection{Was ist Intelligenz?}
\subsubsection*{Turing Test}
\fatsf{Zentrale Frage:} Wann verhält sich ein System intelligent?\\
\fatsf{Annahme:} Eine Entität gilt als intelligent, wenn sie allein durch ihr beobachtbares Verhalten nicht von einer anderen intelligenten Entität zu unterscheiden ist\\
\fatsf{Der Test:}
\begin{enumerate}
    \item Ein menschlicher Fragesteller kommuniziert verdeckt mit zwei Teilnehmnern, A und B. Einer der beiden ist ein Computer, der andere ein Mensch
    \item Gelingt es dem Fragesteller nicht, eindeutig zu bestimmen, welcher der Teilnehmer der Computer ist, gilt der Computer als intelligent und hat den Test bestanden
\end{enumerate}
\fatsf{Relevanz und Kritik:}
\begin{itemize}
    \item Der Turing-Test ist nach wie vor relevant, könnte aber die falsche Fragestellung sein
    \item Er erfordert die Zusammenarbeit zentraler KI-Komponente wie Wissen, logisches Denken, Sprachverständnis und Lernen
    \item Allerdings ist der Test schwer reproduzierbar, nicht konstruktiv und kaum mathematisch analysierbar
\end{itemize}
\fatsf{Weiterführende Frage:} Bedeutet menschenähnliches Verhalten zwangsläufig Intelligenz?
\subsubsection*{Das Chinese-Room-Argument}
\fatsf{Zentrale Frage:} Ist Intelligenz dasselbe wie intelligentes Verhalten?\\
\fatsf{Annahme:} Auch wenn eine Maschine sich intelligent verhält, bedeutet das nicht zwangsläufig, dass sie tatsächlich intelligent ist\\
\fatsf{Der Test:}
\begin{enumerate}
    \item Eine Person, die kein Chinesisch versteht, wird in einen abgeschlossenen Raum gesperrt
    \item Eine andere Person außerhalb des Raums kann schriftliche Notizen auf Chinesisch in den Raum geben, hat aber keine weitere Interaktionsmöglichkeit mit der Person im Raum
    \item Die Person im Raum verfügt über detaillierte Anweisungen, wie sie jede eingehende Nachricht beantworten kann, ohne die Bedeutung zu übersetzen oder zu verstehen
\end{enumerate}
\fatsf{Weiterführende Fragen:}
\begin{itemize}
    \item Ist die Person im Raum intelligent?
    \item Ist ein selbstfahrendes Auto intelligent?
\end{itemize}
\subsubsection*{Müssen wir wissen, was ein intelligentes System ist?}
Die Frage, wie ein intelligentes Wesen definiert wird, ist komplex und seit Langem umstritten\\
Allerdings ist es \fatsf{unwahrscheinlich}, dass diese Definition einen größeren Einfluss auf die Praxis der KI-Forschung hat - genauso wie die Wissenschaftsphilosophie in der Regel nur begrenzten Einfluss auf die tatsächliche wissenschaftliche Praxis hat
\subsection{Merkmale der Künstlichen Intelligenz}
\subsubsection*{Eine KI, sie alle zu beherrschen?}
\fatsf{Allgemeine vs. Spezialisierte KI}
\begin{itemize}
    \item Eine \fatsf{allgemeine KI} (starke KI) ist so konzipiert, dass sie jede intellektuelle Aufgabe bewältigen kann
    \item Eine \fatsf{spezialisierte KI} (schwache KI) ist auf eine oder mehrere festgelegte Aufgaben beschränkt
    \item Diese entspricht der Unterscheidung zwischen ein \fatsf{intelligent sein} und \fatsf{sich intelligent verhalten}
    \item Während starke KI ein langfristiges Forschungsziel ist, verwenden wir derzeit hauptsächlich spezialisierte KI
\end{itemize}
\fatsf{Eigenschaften einer KI}
Eine KI sollte ...
\begin{itemize}
    \item \fatsf{Anpassungsfähig sein}
    \begin{itemize}
        \item Fähigkeit, ihre Leistung durch Lernen aus Erfahrung zu verbessern
    \end{itemize}
    \item \fatsf{Autonom sein}
    \begin{itemize}
        \item Fähigkeit, Aufgaben in verschiedenen Umgebungen ohne ständige menschliche Anleitung auszuführen
    \end{itemize}
    \item \fatsf{Rational handeln}
    \begin{itemize}
        \item \fatsf{Gesetz des Denkens:}
        \begin{itemize}
            \item Grundlage des logischen Denkens und der Frage, welche Argumentations- und Denkprozesse \enquote{korrekt} sind
        \end{itemize}
        \item \fatsf{Rationales Verhalten:}
        \begin{itemize}
            \item Frage nach dem \enquote{richtigen Handeln}
            \item In rational handelnden Systemen bedeutet \enquote{richtig} das, was den Erfolg unter gegebenen Informationen maximiert
        \end{itemize}
    \end{itemize}
\end{itemize}
\fatsf{Vorteile rationalen Verhaltens gegenüber dem Gesetz des Denkens:}
\begin{itemize}
    \item \fatsf{Allgemeiner:} In vielen Situationen gibt es keine nachweisbar richtige Handlung
    \item \fatsf{Besser definierbar und optimierbar:} Rationalität kann mathematisch formuliert und verbessert werden
\end{itemize}
\subsubsection*{Was ist KI?}
Es gibt unterschiedliche Definitionen, je nach gewähltem Kriterium. Zwei zentrale Dimensionen:
\begin{enumerate}
    \item Denken/Schlussfolgern vs. Verhalten/Handeln
    \item Erfolg nach menschlichen Maßstäben vs. Erfolg nach einem idealen Konzept von Intelligenz (Rationalität)
\end{enumerate}
\begin{figure}[h]
    \centering
    \includegraphics[width=1\textwidth]{KI-Systeme Images/So what is AI?.jpeg}
\end{figure}
\subsection{Kognitionswissenschaft}
\subsubsection*{Systeme, die wie Menschen denken}
\fatsf{Wie denken Menschen?}
Die Theorie basiert auf wissenschaftlichen Modellen innerer Gehirnaktivitäten \fatsf{(kognitive Modelle):}
\begin{itemize}
    \item Vorhersage und Überprüfung menschlichen Verhaltens
    \item Identifikation kognitiver Prozesse durch neurologisches Daten
\end{itemize}
Die Kognitionswissenschaft verbindet \fatsf{computationale Modelle aus der KI} mit \fatsf{eperimentellen Technikend er Psychologie}, um präzise und testbare Theorien über den Geist zu entwickeln\\
Oft wird kognitionswissenschaft als eigenständig betrachtet, doch das ist falsch:\\
Wie an der \fatsf{TU Darmstadt} betrachtet, sind KI und Kognitionswissenschaft \fatsf{Zwillingsdisziplinen}\\
\fatsf{Kognitive Neurowissenschaften:} Untersuchung der Funktionsweise des Gehirns auf neuronaler Ebene
\subsection{Grundlagen des Künstlichen Intelligenz}
Verschiedene Disziplinen haben zur Entwicklung der KI beigetragen:
\begin{itemize}
    \item \fatsf{Philosophie:} Logik, Schlussfolgerungen, Geist als physikalisches System, Grundlagen des Lernens, Sprache und Rationalität
    \item \fatsf{Mathematik:} Formale Repräsentation und Beweisverfahren, Berechenbarkeit, (Un-)Entscheidbarkeit, Komplexitätstheorie, Wahrscheinlichkeitstheorie
    \item \fatsf{Psychologie:} Anpassung, Wahrnehmung, Motorik
    \item \fatsf{Ökonomie:} Theorie rationaler Entscheidungen, Spieltheorie
    \item \fatsf{Linguistik:} Wissensrepräsentation, Grammatik
    \item \fatsf{Neurowissenschaften:} Physikalische Grundlage mentaler Prozesse
    \item \fatsf{Regelungstechnik:} Homöostatische Systeme, Stbilität, optimaler Agentenentwurf
\end{itemize}
\subsection{Verwandte Bereiche und Taxonomie der KI}
\subsubsection*{Taxonomie of der KI}
\begin{figure}[h]
    \centering
    \includegraphics[width=1\textwidth]{KI-Systeme Images/Taxonomie of AI.jpeg}
\end{figure}
\subsubsection*{Teilbereiche der KI}
\begin{itemize}
    \item Maschinelles Lernen
    \item Deep Learning
    \item Suche \& Optimierung
    \item Robotik
    \item Verarbeitung natürlicher Sprache (NLP)
    \item Computer Vision
    \item Kognitionswissenschaft
    \item KI-Systeme
    \item Datenwissenschaft
    \item ...
\end{itemize}
\subsection{Angewandte KI in modernen Bereichen}
\begin{itemize}
    \item \fatsf{Autonome Planung und Terminierung}
    \item \fatsf{Spiele} (Planung, Schlussfolgerung, Verständnis ...)
    \item \fatsf{Verarbeitung natürlicher Sprache} (Textgenerierung, Übersetzung, Spracherkennung ...)
    \item \fatsf{Robotik} (Autonomes Fahren ...)
    \item \fatsf{Wissenschaftliche Entdeckungen} (KI findet neue Quasare, beweist Theoreme ...)
    \item \fatsf{Computer Vision} (Autonome Fahrzeuge, Klassifikation, Segmentierung ...)
    \item ...
\end{itemize}
\subsection{Was KI kann und nicht kann}
\enquote{KI ist schwieriger, als man denkt}
\begin{itemize}
    \item Aktuelle KI ist meist auf \fatsf{spezifische} Probleme beschränkt
    \item KI ist \fatsf{nicht} in jeder Aufgabe \fatsf{übermenschlich}
    \item KI-Modelle sind \fatsf{nicht frei von Verzerrungen (Bias)}
    \item Fundamentale Unterschiede in der \fatsf{Wahrnehmung der Welt} zwischen KI und Menschen
\end{itemize}